%% ---------------------------------------------------------------------------
\section{Lines and the Half-Line}\label{sec:lines}

\subsection{Rules, Lines and Synchronization}

Let $S$ be a finite set of \emph{states} that contains three distinct states:
the \emph{quiescent} state $\Lq$, the \emph{general} $\Gq$ and the
\emph{firing} state $\Fq$. The other states are called \emph{auxiliary}; the
rules of this paper use $\Aq$, $\Bq$ and, for Mazoyer's solution, $\Cq$. The
\emph{working} states are $W:=S\setminus\{\Fq\}$, and $\ast\notin S$ denotes the
border. A \emph{neighbourhood} is a triple $(x,y,z)\in(W\cup\{\ast\})\times
W\times(W\cup\{\ast\})$ with $(x,z)\ne(\ast,\ast)$, and a \emph{rule} is a map
$\delta$ from the set of neighbourhoods to $S$ such that
\begin{equation}\label{eq:quiescence}
\delta(\Lq,\Lq,\Lq)=\delta(\Lq,\Lq,\ast)=\Lq .
\end{equation}
If $|S|=k$, there are $k(k-1)k-(k-1)=(k-1)^2(k+1)$ neighbourhoods: $45$ for
four states, $96$ for five and $175$ for six.

The \emph{line of length $n\ge2$} is the partial map
$C_n:\mathbb N\times\{1,\dots,n\}\to S$ given by $C_n(0,1)=\Gq$, $C_n(0,i)=\Lq$
for $2\le i\le n$, and
\[
C_n(t+1,i)=\delta\bigl(C_n(t,i-1),\,C_n(t,i),\,C_n(t,i+1)\bigr)
\]
if the three arguments are defined and none of them is $\Fq$; otherwise
$C_n(t+1,i)$ is undefined. Here $C_n(t,0)=C_n(t,n+1)=\ast$ stand for the
borders. Thus a cell is defined as long as nothing in its neighbourhood has
fired. The rule \emph{synchronizes} the line of length $n$ at time $T$ if
$C_n(T,i)=\Fq$ for $1\le i\le n$ and $C_n(t,i)\in W$ for all $t<T$ and $1\le
i\le n$. A \emph{minimal-time solution} is a rule that synchronizes the line of
length $n$ at time $2n-2$ for every $n\ge2$; it has $k$ states if $|S|=k$.

The \emph{half-line} $C_\infty:\mathbb N\times\{1,2,\dots\}\to S$ is defined in
the same way, with $C_\infty(0,1)=\Gq$, $C_\infty(0,i)=\Lq$ for $i\ge2$, and a
border on the left only. We draw space-time diagrams with time running
downwards and the cells from left to right, as in \cref{fig:line}. The
\emph{front} is the diagonal $t=i-1$, and the \emph{anti-diagonal} $t+i=m$ is
the set of cells $(t,i)$ with $t+i=m$.

\subsection{Agreement}

Information travels at most one cell per step, so the cells that the general
cannot yet have influenced stay quiescent.

\begin{lemma}[light cone]\label{lem:lightcone}
For every rule, $C_n(t,i)=C_\infty(t,i)=\Lq$ whenever $t<i-1$ (and $i\le n$
for $C_n$).
\end{lemma}

\begin{proof}
By induction on $t$. For $t=0$ this is the initial condition. For $t\ge1$, the
cells $(t-1,i')$ with $i'\in\{i-1,i,i+1\}$ satisfy $t-1<i'-1$, so each of them
is $\Lq$ or the right border, and \eqref{eq:quiescence} gives
$C_n(t,i)=C_\infty(t,i)=\Lq$.
\end{proof}

The right border of the line of length $n$ is first felt by the cell $n$ at
time $n-1$, when the front reaches it, and its influence travels back to the
left at speed at most one. Before this echo arrives, the line cannot be told
apart from the half-line.

\begin{lemma}[agreement]\label{lem:agree}
For every rule, every $n\ge2$, $1\le i\le n$ and $t\ge0$ with $t+i\le 2n-2$,
the values $C_n(t,i)$ and $C_\infty(t,i)$ are both undefined or both defined
and equal.
\end{lemma}

\begin{proof}
By induction on $t$; the case $t=0$ is clear. Let $t\ge1$. If $i=n$, then
$t\le n-2<i-1$, and both values are $\Lq$ by \cref{lem:lightcone}. If $i<n$,
the cells $(t-1,j)$, $j\in\{i-1,i,i+1\}$, satisfy $(t-1)+j\le 2n-2$ and $j\le
n$, so by induction they are both undefined or equal in $C_n$ and $C_\infty$
(for $j=0$ both are the border), and the recursion gives the claim.
\end{proof}

So the border can only matter in the \emph{reflected triangle}
\[
R_n:=\{(t,i):\ 1\le i\le n,\ t+i\ge 2n-1,\ t\le 2n-2\},
\]
which is bounded by the anti-diagonal $t+i=2n-1$, the right border and the
firing row $t=2n-2$ (\cref{fig:line}). The agreement lemma already shows why
time $2n-2$ is optimal.

\begin{remark}\label{rem:lowerbound}
Suppose that a rule synchronizes every line, the line of length $n$ at some
time $T\le 2n-3$. Since $T+1\le 2n-2$, \cref{lem:agree} gives
$C_\infty(T,1)=C_n(T,1)=\Fq$, and then $C_{n'}(T,1)=\Fq$ for every $n'\ge n$,
by \cref{lem:agree} again. For $n'>T+1$ the cell $n'$ is still quiescent at
time $T$ by \cref{lem:lightcone}, so the line of length $n'$ has a cell in
state $\Fq$ at time $T$ although not all of its cells are in state $\Fq$; it
can then never be synchronized, a contradiction (\cref{fig:lowerbound}).
\end{remark}

\begin{figure}[t]
\centering
%% Figure: why no line is synchronized before time 2n-2 (Remark 2.3). Schematic, with n = 8, a
%% hypothetical firing time T = 11 <= 2n-3 and n' = 13 > T+1; cell (t,i) is [i-1,i] x [-t-1,-t].
\begin{tikzpicture}[x=3.1mm,y=3.1mm]
  % (a) the line of length n = 8, synchronized at time T = 11
  \begin{scope}
    \foreach \t in {0,...,11} {\foreach \i in {1,...,8} {%
      \pgfmathtruncatemacro{\s}{\t+\i}\pgfmathtruncatemacro{\lc}{\i-1}%
      \def\cc{rgRefl}%
      \ifnum\s<15 \def\cc{rgAgree}\fi
      \ifnum\t=11 \def\cc{stF}\fi
      \ifnum\t<\lc \def\cc{white}\fi
      \filldraw[fill=\cc,draw=grid!45,line width=0.15pt] (\i-1,-\t) rectangle (\i,-\t-1);}}
    \begin{scope}\clip (0,0) rectangle (8,-12);
      \draw[white,line width=1.8pt] \staircase{14}{7} -- (8,-7);
      \draw[ink,densely dashed,line width=0.8pt] \staircase{14}{7} -- (8,-7);
    \end{scope}
    \draw[ink,line width=0.6pt] (0,0) rectangle (8,-12);
    \node[font=\small,anchor=south] at (4,0.4) {(a) line of length $n$};
    \node[font=\scriptsize,anchor=east] at (-0.3,-0.5) {$t=0$};
    \node[font=\scriptsize,anchor=east] at (-0.3,-11.5) {$t=T$};
    \node[font=\scriptsize,anchor=north,align=center] at (4,-12.4) {synchronized at a\\ time $T\le 2n-3$};
  \end{scope}
  % (b) the half-line up to time T
  \begin{scope}[shift={(12,0)}]
    \foreach \t in {0,...,11} {\foreach \i in {1,...,14} {%
      \pgfmathtruncatemacro{\s}{\t+\i}\pgfmathtruncatemacro{\lc}{\i-1}%
      \def\cc{gNone}%
      \ifnum\s<15 \def\cc{rgAgree}\ifnum\t=11 \def\cc{stF}\fi\fi
      \ifnum\t<\lc \def\cc{white}\fi
      \filldraw[fill=\cc,draw=grid!45,line width=0.15pt] (\i-1,-\t) rectangle (\i,-\t-1);}}
    \begin{scope}\clip (0,0) rectangle (14,-12);
      \draw[white,line width=1.8pt] \staircase{14}{14};
      \draw[ink,densely dashed,line width=0.8pt] \staircase{14}{14};
    \end{scope}
    \draw[ink,line width=0.6pt] (0,0) -- (0,-12) -- (14,-12) (0,0) -- (14,0);
    \node[font=\scriptsize,fill=white,inner sep=1pt,rotate=45] at (9.6,-4.2) {$t+i=2n-2$};
    \node[font=\small,anchor=south] at (7,0.4) {(b) half-line};
    \node[font=\scriptsize,anchor=north,align=center] at (7,-12.4) {$C_\infty(T,1)=C_n(T,1)=\Fq$\\
      (\cref{lem:agree}, as $T+1\le 2n-2$)};
  \end{scope}
  % (c) the line of length n' = 13 > T+1 up to time T
  \begin{scope}[shift={(30,0)}]
    \foreach \t in {0,...,11} {\foreach \i in {1,...,13} {%
      \pgfmathtruncatemacro{\s}{\t+\i}\pgfmathtruncatemacro{\lc}{\i-1}%
      \def\cc{gNone}%
      \ifnum\s<15 \def\cc{rgAgree}\ifnum\t=11 \def\cc{stF}\fi\fi
      \ifnum\t<\lc \def\cc{white}\fi
      \filldraw[fill=\cc,draw=grid!45,line width=0.15pt] (\i-1,-\t) rectangle (\i,-\t-1);}}
    \begin{scope}\clip (0,0) rectangle (13,-12);
      \draw[white,line width=1.8pt] \staircase{14}{13};
      \draw[ink,densely dashed,line width=0.8pt] \staircase{14}{13};
    \end{scope}
    \draw[ink,line width=0.6pt] (0,0) rectangle (13,-12);
    \draw[ink,line width=1.1pt] (12,-11) rectangle (13,-12);
    \node[font=\small,anchor=south] at (6.5,0.4) {(c) line of length $n'>T+1$};
    \node[font=\scriptsize,anchor=north,align=center] at (6.5,-12.4) {$C_{n'}(T,1)=\Fq$, but the cell $n'$\\
      is still quiescent (\cref{lem:lightcone})};
  \end{scope}
  % legend
  \node[anchor=north,font=\scriptsize,align=left] at (21.5,-15.6) {%
    \swatch{white}\,quiescent, $t<i-1$\qquad
    \swatch{rgAgree}\,$t+i\le 2n-2$: the same in (a), (b) and (c) (\cref{lem:agree})\qquad
    \swatch{stF}\,$\Fq$\\[0.4ex]
    \swatch{rgRefl}\,the reflected triangle $R_n$\qquad
    \swatch{gNone}\,the same in (b) and (c) (\cref{lem:agree} for the length $n'$)};
\end{tikzpicture}

\caption{Why no line can be synchronized before time $2n-2$
(\cref{rem:lowerbound}), drawn for $n=8$, $T=11$ and $n'=13$. (a)~If the line
of length $n$ were synchronized at a time $T\le 2n-3$, its cell $(T,1)$ would
satisfy $T+1\le2n-2$ and lie above the dashed staircase, where the line agrees
with the half-line. (b)~Hence the cell $(T,1)$ of the half-line would be in
state $\Fq$. (c)~A line of length $n'>T+1$ agrees with the half-line on every
cell with $t\le T$, because then $t+i\le T+n'\le 2n'-2$. Its cell $1$ would be
in state $\Fq$ at time $T$ while its cell $n'$ is still quiescent, so this line
could never be synchronized.}
\label{fig:lowerbound}
\end{figure}

\subsection{The Return Chain}

The first anti-diagonal of $R_n$ is where the line starts to differ from the
half-line, and in a minimal-time solution it differs there everywhere
(\cref{fig:line}(c) and \cref{fig:return}).

\begin{lemma}[return chain]\label{lem:return}
Let $n\ge2$. If $\delta$ synchronizes the lines of lengths $n$ and $n+1$ at
times $2n-2$ and $2n$, then $C_n(2n-1-i,i)\ne C_\infty(2n-1-i,i)$ for $1\le i\le
n$, and the front cell $f_{n-1}:=C_\infty(n-2,n-1)$ is not $\Lq$.
\end{lemma}

\begin{proof}
For $i=1$ we have $C_n(2n-2,1)=\Fq$, whereas
$C_\infty(2n-2,1)=C_{n+1}(2n-2,1)\ne\Fq$ by \cref{lem:agree} (applied to the
length $n+1$) and because $2n-2<2n$. Suppose that the claim holds for some
$i<n$. The left and middle predecessors of the cell $(2n-1-i,i)$, namely
$(2n-2-i,i-1)$ and $(2n-2-i,i)$, lie on the anti-diagonals $t+i=2n-3$ and
$2n-2$ (for $i=1$ the left one is the border), so they agree in $C_n$ and
$C_\infty$ by \cref{lem:agree}. As the cell itself differs, so do the right
predecessors $(2n-2-i,i+1)$; all these cells are defined, since no cell of the
line of length $n$ fires before time $2n-2$. This is the claim for $i+1$. For
$i=n$ the claim concerns the cell $(n-1,n)$. Its predecessors satisfy
$C_n(n-2,n-1)=C_\infty(n-2,n-1)=f_{n-1}$ by \cref{lem:agree}, and the cells
$(n-2,n)$ and $(n-2,n+1)$ are quiescent by \cref{lem:lightcone}, so the claim
reads $\delta(f_{n-1},\Lq,\ast)\ne\delta(f_{n-1},\Lq,\Lq)$. This is impossible
if $f_{n-1}=\Lq$, by \eqref{eq:quiescence}.
\end{proof}

\begin{figure}[t]
\centering
%% Figure: the induction of Lemma 2.4 along the return chain. Schematic, n = 8; cell (t,i) is
%% [i-1,i] x [-t-1,-t] in panel (a).
\begin{tikzpicture}
  % (a) the line of length n = 8
  \begin{scope}[x=3.4mm,y=3.4mm]
    \foreach \t in {0,...,14} {\foreach \i in {1,...,8} {%
      \pgfmathtruncatemacro{\s}{\t+\i}\pgfmathtruncatemacro{\lc}{\i-1}%
      \def\cc{rgRefl}%
      \ifnum\s<16 \def\cc{rgReturn}\fi
      \ifnum\s<15 \def\cc{rgInput}\fi
      \ifnum\s<13 \def\cc{rgAgree}\ifnum\t<\lc \def\cc{white}\fi\fi
      \ifnum\t=14 \def\cc{stF}\fi
      \filldraw[fill=\cc,draw=grid!45,line width=0.15pt] (\i-1,-\t) rectangle (\i,-\t-1);}}
    \draw[ink,line width=0.6pt] (0,0) rectangle (8,-15);
    % the step of panel (b), for i = 3
    \foreach \x/\y in {1/11,2/11,3/11} \draw[ink,line width=0.9pt,densely dashed] (\x,-\y) rectangle (\x+1,-\y-1);
    \draw[ink,line width=1.3pt] (2,-12) rectangle (3,-13);
    % the front cell f_{n-1} = (n-2, n-1) and the last cell (n-1, n) of the return chain
    \node[font=\scriptsize] at (6.5,-6.5) {$f$};
    \draw[ink,line width=1.3pt] (7,-7) rectangle (8,-8);
    % direction of the induction
    \draw[-{Stealth[length=1.8mm]},ink,line width=0.7pt] (1.9,-15.6) -- (9.0,-8.5);
    \node[font=\scriptsize,rotate=45,anchor=south] at (5.2,-11.9) {induction on $i$};
    \node[font=\small,anchor=south] at (4,0.4) {(a) line of length $n$};
    \foreach \y/\l in {0/{t=0},7/{t=n-1},14/{t=2n-2}}
      \node[font=\scriptsize,anchor=east] at (-0.3,-\y-0.5) {$\l$};
    \foreach \x/\l in {1/1,8/n}
      \node[font=\scriptsize,anchor=north] at (\x-0.5,-15.2) {$\l$};
  \end{scope}
  % (b) the induction step
  \begin{scope}[shift={(4.7,-0.9)},x=9.5mm,y=9.5mm]
    \filldraw[fill=rgInput,draw=ink,line width=0.5pt] (0,0) rectangle (1,-1);
    \filldraw[fill=rgInput,draw=ink,line width=0.5pt] (1,0) rectangle (2,-1);
    \filldraw[fill=rgReturn,draw=ink,line width=0.5pt] (2,0) rectangle (3,-1);
    \filldraw[fill=rgReturn,draw=ink,line width=0.5pt] (1,-1.6) rectangle (2,-2.6);
    \node[font=\small] at (0.5,-0.5) {$=$};
    \node[font=\small] at (1.5,-0.5) {$=$};
    \node[font=\small,text=white] at (2.5,-0.5) {$\neq$};
    \node[font=\small,text=white] at (1.5,-2.1) {$\neq$};
    \foreach \x in {0.5,1.5,2.5} \draw[-{Stealth[length=1.5mm]},ink,line width=0.5pt] (\x,-1.03) -- (1.5,-1.57);
    \node[font=\scriptsize,anchor=south] at (0.5,0.02) {$2n-3$};
    \node[font=\scriptsize,anchor=south] at (1.5,0.02) {$2n-2$};
    \node[font=\scriptsize,anchor=south] at (2.5,0.02) {$2n-1$};
    \node[font=\scriptsize,anchor=east] at (-0.1,0.25) {$t+i$:};
    \node[font=\scriptsize,anchor=west,align=left] at (3.1,-0.5) {cell $i+1$ of the\\ return chain};
    \node[font=\scriptsize,anchor=west,align=left] at (2.1,-2.1) {cell $i$ of the\\ return chain};
    \node[font=\small,anchor=south west] at (-0.9,0.7) {(b) the induction step};
  \end{scope}
  % (c) the last cell of the return chain
  \begin{scope}[shift={(4.7,-4.35)},x=9.5mm,y=9.5mm]
    \foreach \row/\z/\lab in {0/{\ast}/{line of length $n$}, 1.15/{\Lq}/{half-line}} {
      \filldraw[fill=rgInput,draw=ink,line width=0.5pt] (0,-\row) rectangle (1,-\row-1);
      \filldraw[fill=white,draw=ink,line width=0.5pt] (1,-\row) rectangle (2,-\row-1);
      \filldraw[fill=white,draw=ink,line width=0.5pt] (2,-\row) rectangle (3,-\row-1);
      \node[font=\small] at (0.5,-\row-0.5) {$f$};
      \node[font=\small] at (1.5,-\row-0.5) {$\Lq$};
      \node[font=\small] at (2.5,-\row-0.5) {$\z$};
      \node[font=\scriptsize,anchor=west] at (3.1,-\row-0.5) {\lab};
    }
    \node[font=\small,anchor=south west] at (-0.9,0.25) {(c) the last step};
  \end{scope}
  % legend
  \node[anchor=north,font=\scriptsize,align=left] at (4.9,-6.75) {%
    \swatch{white}\,quiescent\quad \swatch{rgAgree}\,$t+i\le 2n-4$\quad
    \swatch{rgInput}\,$t+i=2n-3$ and $2n-2$\quad \swatch{rgReturn}\,$t+i=2n-1$\quad
    \swatch{rgRefl}\,rest of $R_n$\quad \swatch{stF}\,$t=2n-2$};
\end{tikzpicture}

\caption{The proof of \cref{lem:return}, drawn for $n=8$. (a)~The line of
length $n$. On the cells with $t+i\le2n-2$, including the two anti-diagonals
$t+i=2n-3$ and $t+i=2n-2$, it agrees with the half-line (\cref{lem:agree}). The
return chain $t+i=2n-1$ starts at the cell $(2n-2,1)$, which is $\Fq$ in the
line but not in the half-line, and the induction climbs the chain to the cell
$(n-1,n)$. (b)~The induction step, outlined in (a) for $i=3$: the left and
middle predecessors of the cell $i$ of the return chain have the same state in
the line and in the half-line ($=$); since the cell differs ($\neq$), so does
its right predecessor, the cell $i+1$ of the return chain. (c)~The last cell
$(n-1,n)$ of the return chain, outlined in (a), is computed from $f$, $\Lq$,
$\ast$ in the line and from $f$, $\Lq$, $\Lq$ in the half-line, where
$f=f_{n-1}$ is the front cell; by \eqref{eq:quiescence} the results can differ
only if $f\neq\Lq$.}
\label{fig:return}
\end{figure}

We call the anti-diagonal $t+i=2n-1$ the \emph{return chain} of the line of
length $n$. Finally, a minimal-time solution never fires on the half-line:
given a cell $(t,i)$, take $n\ge t+i+1$; then $C_\infty(t,i)=C_n(t,i)$ by
\cref{lem:agree}, and $t<2n-2$, so
\begin{equation}\label{eq:neverfire}
C_\infty(t,i)\in W\qquad\text{for all }t\ge0\text{ and }i\ge1 .
\end{equation}
